% Można zmienić na angielski wybierając drugą linię
\documentclass[11pt]{aghdpl}
% \documentclass[en,11pt]{aghdpl}

% Lista wszystkich języków stanowiących języki pozycji bibliograficznych użytych w pracy.
% (Zgodnie z zasadami tworzenia bibliografii każda pozycja powinna zostać utworzona zgodnie z zasadami języka, w którym dana publikacja została napisana.)
\usepackage[english,polish]{babel}

% Użyj polskiego łamania wyrazów (zamiast domyślnego angielskiego).
% Można zmienić na angielski
\usepackage{polski}

\usepackage[utf8]{inputenc}

% dodatkowe pakiety

\usepackage{mathtools}
\usepackage{amsfonts}
\usepackage{amsmath}
\usepackage{amsthm}


% Paczka wyłącza pokazywanie numeru strony oraz nagłówków na pustych stronach
% Można wyłączyć
\usepackage{emptypage}

% Dodatkowe kolory:
\usepackage{colortbl}
\definecolor{redi}{RGB}{255,38,0}
\definecolor{redii}{RGB}{200,50,30}
\definecolor{yellowi}{RGB}{255,251,0}
\definecolor{bluei}{RGB}{0,150,255}
\definecolor{blueii}{RGB}{135,247,210}
\definecolor{blueiii}{RGB}{91,205,250}
\definecolor{blueiv}{RGB}{115,244,253}
\definecolor{bluev}{RGB}{1,58,215}
\definecolor{orangei}{RGB}{240,143,50}
\definecolor{yellowii}{RGB}{222,247,100}
\definecolor{greeni}{RGB}{166,247,166}
\definecolor{gray}{RGB}{200, 200, 200}
\definecolor{lgray}{RGB}{240, 240, 240}

\usepackage{siunitx}

% Ustawienia marginesów kolumn w array
% Można usunąć/zmodyfikować
\usepackage{array}
\newcolumntype{0}{>{\centering\arraybackslash}m{0.5cm}}
\newcolumntype{1}{>{\centering\arraybackslash}m{1cm}}
\newcolumntype{Q}{>{\centering\arraybackslash}m{1.5cm}}
\newcolumntype{2}{>{\centering\arraybackslash}m{2cm}}
\newcolumntype{W}{>{\centering\arraybackslash}m{2.5cm}}
\newcolumntype{3}{>{\centering\arraybackslash}m{3cm}}
\newcolumntype{4}{>{\centering\arraybackslash}m{4cm}}
\newcolumntype{5}{>{\centering\arraybackslash}m{5cm}}
\newcolumntype{6}{>{\centering\arraybackslash}m{6cm}}
\newcolumntype{7}{>{\centering\arraybackslash}m{7cm}}
\newcolumntype{8}{>{\centering\arraybackslash}m{8cm}}

% Dzięki temu będzie się dało kopiować tekst który ma polskie literki
% see: https://tex.stackexchange.com/questions/57915/cannot-copy-letters-with-diacritics-from-pdflatex-pdf
\usepackage{lmodern}
\usepackage{listings}

% Polskie literki w listingach:
\lstset{
    literate=
    {ą}{{\k{a}}}1 {Ą}{{\k{A}}}1 
    {ł}{{\l{}}}1 {Ł}{{\L{}}}1 
    {ń}{{\'n}}1 {Ń}{{\'N}}1 
    {ę}{{\k{e}}}1 {Ę}{{\k{E}}}1 
    {ś}{{\'s}}1 {Ś}{{\'S}}1 
    {ż}{{\.z}}1 {Ż}{{\.Z}}1 
    {ó}{{\'o}}1 {Ó}{{\'O}}1 
    {ź}{{\'z}}1 {Ź}{{\'Z}}1 
    {ć}{{\'c}}1 {Ć}{{\'C}}1
}

% Kolorowanie składni języka C++
% TODO: Można usunąć jeśli się nie potrzebuje
\definecolor{listinggray}{gray}{0.9}
\definecolor{lbcolor}{rgb}{0.9,0.9,0.9}
\lstset{
    language=C++,
    basicstyle=\ttfamily,
    breaklines=true,
    breakatwhitespace=true, %%%
    basicstyle=\footnotesize,
    %upquote=true,
    aboveskip={0.75\baselineskip},
    columns=fixed,
    showstringspaces=false,
    %extendedchars=false,
    breaklines=true,
    prebreak = \raisebox{0ex}[0ex][0ex]{\ensuremath{\hookleftarrow}},
    %frame=single,
    numbers=left,
    showtabs=false,
    showspaces=false,
    showstringspaces=false,
    %backgroundcolor=\color{lbcolor},
    identifierstyle=\ttfamily,
    keywordstyle=\color[rgb]{0,0,1},
    commentstyle=\color[rgb]{0.4,0.4,0.4},
    stringstyle=\color[rgb]{0.627,0.126,0.941},
    numberstyle={\footnotesize\color[rgb]{0.205, 0.142, 0.73}}
}

\usepackage[hidelinks]{hyperref}

% --- < bibliografia > ---
%
% Dla TeXstudio warto przeczytać poniższe!
% UWAGA: Żeby bibliografia działała gdy używamy TeXstudio, to należy zmodyfikować quick build w opcjach, tak aby wykonywane były polecenia:
% 	PdfLaTeX
% 	BibTeX
% 	PdfLaTeX
% 	PdfLaTeX
%
% Więcej informacji na:
% 	https://tex.stackexchange.com/a/216325
%

\usepackage[
	backend=bibtex,
	style=numeric,
	sorting=none,
	% Zastosuj styl wpisu bibliograficznego właściwy językowi publikacji.
	language=autobib,
	autolang=other,
	% Zapisuj datę dostępu do strony WWW w formacie RRRR-MM-DD.
	urldate=iso8601,
	% Nie dodawaj numerów stron, na których występuje cytowanie.
	backref=false,
	% Podawaj ISBN.
	isbn=true,
	% Nie podawaj URL-i, o ile nie jest to konieczne.
	url=false,
	% Ustawienia związane z polskimi normami dla bibliografii.
	maxbibnames=3,
]{biblatex}

\usepackage{csquotes}
% Ponieważ `csquotes` nie posiada polskiego stylu, można skorzystać z mocno zbliżonego stylu chorwackiego.
\DeclareQuoteAlias{croatian}{polish}

% Bibliografia musi być uzupełniona w tym pliku:
\addbibresource{bibliografia.bib}

% Nie wyświetlaj wybranych pól.
%\AtEveryBibitem{\clearfield{note}}


% ------------------------
% --- < listingi > ---

% Użyj czcionki kroju Courier.
\usepackage{courier}

\usepackage{listings}
\lstloadlanguages{TeX}

% Polskie literki:
\lstset{
	literate={ą}{{\k{a}}}1
           {ć}{{\'c}}1
           {ę}{{\k{e}}}1
           {ó}{{\'o}}1
           {ń}{{\'n}}1
           {ł}{{\l{}}}1
           {ś}{{\'s}}1
           {ź}{{\'z}}1
           {ż}{{\.z}}1
           {Ą}{{\k{A}}}1
           {Ć}{{\'C}}1
           {Ę}{{\k{E}}}1
           {Ó}{{\'O}}1
           {Ń}{{\'N}}1
           {Ł}{{\L{}}}1
           {Ś}{{\'S}}1
           {Ź}{{\'Z}}1
           {Ż}{{\.Z}}1,
	basicstyle=\footnotesize\ttfamily,
}

% ------------------------
% Określamy nazwy 'table' oraz 'figure':
\AtBeginDocument{
	\renewcommand{\tablename}{Tabela}
	\renewcommand{\figurename}{Rys.}
}

% ------------------------
% --- < tabele > ---

% Dodatkowe paczki do tabel, można usunąć
\usepackage{array}
\usepackage{tabularx}
\usepackage{multirow}
\usepackage{booktabs}
\usepackage{makecell}
\usepackage[flushleft]{threeparttable}

% defines the X column to use m (\parbox[c]) instead of p (`parbox[t]`)
\newcolumntype{C}[1]{>{\hsize=#1\hsize\centering\arraybackslash}X}


%---------------------------------------------------------------------------

% TODO: Poniższe mogą się wyświetlać lub nie w zależności od wariantów wyżej
% Np. języka lub innej konfiguracji
% Prawdopodobnie niektóre z poniższych nie są nigdy używane.
% Jeśli o tym wiesz, to warto rozwinąć ten szablon i udostępnić go dalej ;)

\author{Karol Kocierz}
\shortauthor{K. Kocierz}

\titlePL{Przewidywanie przestępstw na podstawie danych historycznych i aktualnych}
% \titleEN{TitleEN}

% Skrócone wersje tytułów po polsku/angielsku
% wyświetlane np. gdy normalne są bardzo długie
% \shorttitlePL{Short title PL}
% \shorttitleEN{Short title EN}

% Rodzaj pracy
\thesistype{Praca dyplomowa inżynierska}
%\thesistype{Master of Science Thesis}

% Promotor/opiekun pracy
\supervisor{Dr. Piotr Kowalski}

% Kierunek studiów
\degreeprogramme{Informatyka Stosowana}

% TODO: DATE
\date{2021}

% TODO: Można zakomentować jeśli nie chcemy wyświetlać katedry
% TODO: Nie jest pewne, czy to jest gdzieś wykorzystywane
%\department{TODO: Katedra ~}

% Wydział
\faculty{Wydział Fizyki i Informatyki Stosowanej}

% TODO: Nie jest pewne czy jest wykorzystywane
%\acknowledgements{Serdecznie dziękuję \dots tu ciąg dalszych podziękowań np. dla promotora, żony, sąsiada itp.}


\setlength{\cftsecnumwidth}{10mm}

%---------------------------------------------------------------------------
\setcounter{secnumdepth}{4}

\begin{document}


% =====  STRONA TYTULOWA PRACY INŻYNIERSKIEJ ====
% ostatnia modyfikacja: 2011/03/09, K. Malarz

\thispagestyle{empty}

%% ------------------------ NAGLOWEK STRONY ---------------------------------
\includegraphics[height=37.5mm]{images/agh_nzw_a_pl_1w_wbr_cmyk}\\
\rule{30mm}{0pt}
{\Large\textsf{Wydział Fizyki i Informatyki Stosowanej}}\\
\rule{\textwidth}{3pt}\\
\rule[2.75ex]
{\textwidth}{1pt}%\\
\vspace{1em}\rule{30mm}{0pt}
\begin{center}
    {\bf\huge\textsf{Praca inżynierska}}\\
    \vspace{13ex}
    % --------------------------- IMIE I NAZWISKO -------------------------------
    {\bf\LARGE\textsf{Karol Kocierz}}\\
    \vspace{3ex}
    {\sf \normalsize kierunek studiów:} {\bf\normalsize\textsf{Informatyka Stosowana}}\\
    \vspace{15ex}
    %% ------------------------ TYTUL PRACY --------------------------------------
    {\bf\Huge\textsf{Przewidywanie przestępstw na podstawie danych historycznych i aktualnych}}\\
    \vspace{14ex}
    %% ------------------------ OPIEKUN PRACY ------------------------------------
    {\sf \LARGE Opiekun:} {\bf\LARGE\textsf{Dr. Piotr Kowalski}}\\
    \vspace{22ex}
    \textsf{\bf\Large\textsf{Kraków, Styczeń 2021}}
\end{center}
%% =====  STRONA TYTUŁOWA PRACY INŻYNIERSKIEJ  ====

\newpage

%% =====  TYŁ STRONY TYTUŁOWEJ PRACY INŻYNIERSKIEJ  ====
{\sf Oświadczam, świadomy odpowiedzialności karnej za~poświadczenie nieprawdy, że~niniejszą pracę dyplomową wykonałem osobiście i~samodzielnie i~nie korzystałem ze~źródeł innych niż wymienione w~pracy.}

\vspace{14ex}

\begin{center}
    \begin{tabular}{lr}
        ~~~~~~~~~~~~~~~~~~~~~~~~~~~~~~~~~~~~~~~~~~~~~~~~~~~~~~~~~~~~~~~~~ &
        ................................................................. \\
        ~ & {\sf (czytelny podpis)} \\
    \end{tabular}
\end{center}

% =====  TYL STRONY TYTULOWEJ PRACY INŻYNIERSKIEJ  ====
\newpage
%\titlepages

% Ponowne zdefiniowanie stylu `plain`, aby usunąć numer strony z pierwszej strony spisu treści i poszczególnych rozdziałów.
\fancypagestyle{plain}
{
	% Usuń nagłówek i stopkę
	\fancyhf{}
	% Usuń linie.
	\renewcommand{\headrulewidth}{0pt}
	\renewcommand{\footrulewidth}{0pt}
}

\setcounter{tocdepth}{2}
\tableofcontents
\clearpage

\paragraph{}
Coraz powszechniejszy dostęp do dużych ilości danych z wszelkich dziedzin pozwala na ekstensywne badanie rzeczywistych zjawisk i świat przestępczy nie jest od tej reguły wyjątkiem. Wiele rządów, komend policji czy innych organizacji udostępnia informacje o popełnianych na danym obszarze przestępstwach które nieżadko bywają bardzo dokładne i rozległe. Tego typu bazę danych upublicznia urząd miasta Baltimore MD z miedzy innymi takimi danymi jak:
\begin{itemize}
  \item Typ przestępstwa
  \item Data i czas przestępstwa
  \item Współrzędne w jakich dokonało się przestępstwo
\end{itemize}
\paragraph{}
W poniższej pracy przeprowadzono analizę tych danych pod kontem przewidywania przestępstw ze względu na datę, godzinę i lokalizację.
\section{Klasteryzacja K-means}
\subsection{Klasteryzacja}
\paragraph{}
Klasteryzacją nazywany jest proces łączenia danych w zestawy tak żeby punkty wewnątrz klastra były jak najbardziej do siebie zbliżone względem wybranego zestawu cech natomiast punkty należące do różnych klastrów mają się od siebie jak najbardziej różnić. Metoda ta jest używana do analizy struktury danych w celu wyodrębnienia danej ilości grup.
\subsection{K-means}
\paragraph{}
Metoda k-means jest uważana za najbardziej powszechnie używaną przy klasteryzacji do czego przyczynia się niski stopień skomplikowania. K-means należy do algorytmów maksymalizacji oczekiwań, to jest algorytmów które składają się z dwóch części: oczekiwań (oznaczone poniżej jako O) i maksymalizacji (oznaczone jako M). Algorytm metody wygląda w sposób następujący:
\begin{enumerate}
	\item Ustalenie ilości szukanych klastrów
	\item Losowa inicjalizacja centroidów
	\item Powtarzane aż klastry przestaną się zmieniać:
	\begin{enumerate}
		\item Przypisanie każdego punktu do klastra którego centroid znajduje się najbliżej niego [O]\\
		\begin{equation}
			l^{(i)} := \arg \underset{j}{\min} \left \| x^{(i)} - c_{j} \right \|^{2}
		\end{equation}
		\item Ustawienie każdego centroidu w punkcie centralnym względem wszystkich punktów jego klastra [M]\\
		\begin{equation}
			c_{j} := \frac{\sum_{i=1}^{m}1\left \{ l^{(i)}=j \right \}x^{(i)}}{\sum_{i=1}^{m}1\left \{ l^{(i)}=j \right \}}
		\end{equation}
	\end{enumerate}
\end{enumerate}
\section{Indeks Calinskiego-Harabasza}
\paragraph{}
Przy wykonywaniu klasteryzacji warto móc sprawdzić jaka metoda klasteryzacji z jakimi parametrami daje najlepsze efekty. W tym celu przydatne są rozmaite indeksy, funkcje których celem jest ocena wyniku klasteryzacji. 
\paragraph{}
Jednym z takich indeksów jest indeks Calinskiego-Harabasza. Który dokonuje analizy klasteryzacji porównując podobieństwa między punktami z tego samego klastra do podobieństw między klastrami. Im większe podobieństwo punktów w klastrze i mniejsze podobieństwo między klastrami tym ocena wyniku jest wyższa. Funkcja opisująca indeks jest następująca:
\begin{equation}
	s = \frac{\mathrm{tr}(B_k)}{\mathrm{tr}(W_k)} \times \frac{n_E - k}{k - 1}
\end{equation}
\paragraph{}
Gdzie $\mathrm{tr}(B_k)$ to ślad macierzy kowariancji między klastrami a $\mathrm{tr}(W_k)$ to ślad macierzy kowariancji wewnątrz klastrów. Obie macierze definiowane są następująco:
\begin{equation}
	W_k = \sum_{q=1}^k \sum_{x \in C_q} (x - c_q) (x - c_q)^T
\end{equation}
\begin{equation}
	B_k = \sum_{q=1}^k n_q (c_q - c_E) (c_q - c_E)^T
\end{equation}
\paragraph{}
Wartości w równaniach 2.3 do 2.5 oznaczają: $n_E$ - liczebność danych, $k$ - ilość klastrów, $C_q$ - liczebność klastra $q$, $c_q$ - środek klastra $q$, $c_E$ - środek całego zestawu danych, $n_q$ - liczebność klastra $q$.
\section{Estymator jądrowy gęstości}
\paragraph{}
Estymator jądrowy gęstości lub KDE(z ang. Kernel Density Estimator) jest nieparametrycznym estymatorem który w celu wyestymowania rozkładu wartości wykorzystuje funkcję dodatnią $K$ (tzw. Jądro). Jądrem często są takie funkcje jak gęstość rozkładu normalnego, trójkątnego czy tak zwane jądro Epanecznikowa które opisuje wzór:
\begin{equation}
	K(x) = \dfrac{3}{4}\left ( 1 - x^2 \right ), gdzie \left | x \right | \leq 1
\end{equation}
Natomiast sam estymator jądrowy opisany jest następującym wzorem:
\begin{equation}
	\hat{f}(x) = \dfrac{1}{mh^n}\sum_{i=1}^{m} K\left ( \dfrac{x - x_i}{h}  \right )
\end{equation}
\paragraph{}
W powyższym równaniu $h$ jest parametrem wygładzania i ma wpływ na to jak znaczącą rolę w wyestymowaniu rozkładu w danym punkcie mają punkty znajdujące się w okół niego. Im jego wartość jest większa tym większe jest wygładzenie wynikowego rozkładu.
\paragraph{}
KDE sprawdza się świetnie gdy chcemy dużą ilość skumulowanych danych zwizualizować tak żeby móc wskazać duże nagromadzenia podobnych wartości.
\paragraph{}
Dane poddane badaniu w tej pracy pochodzą z bazy danych urzędu miasta Baltimore a dokładniej jest to rejestr wszystkich zgłoszonych policji tego miasta przestępstw od roku **** do 2020. Rekordy tej bazy zawierają poniższe atrybuty:


\printbibliography

\end{document}
